\begin{statementbox}
\textbf{\textcolor{CStatement}{条件与符号}}

抛物线 $y^2=2px$，$p>0$；焦点 $F(p/2,0)$，准线 $x=-p/2$。
过 $F$ 的直线与抛物线相交于两个不同的有限点 $A,B$。
记两段正长度 $r=AF$、$s=BF$；焦点在线段 $AB$ 内，故 $AB=r+s$。

\par\medskip
\textbf{\textcolor{CStatement}{核心结论：倒数之和不随弦的方向改变}}
\[
 \boxed{\frac1r+\frac1s=\frac2p},\qquad
 \boxed{2rs=p(r+s)}.
\]
这里 $p$ 是焦点到准线的距离；$2p$ 是方程中 $x$ 项的系数，也是通径长。

\par\medskip
\textbf{\textcolor{CStatement}{三条直接推论}}
\[
 AB=r+s=\frac{2rs}{p}\ge2p,
 \qquad AB=2p\iff r=s=p\quad\text{（通径）}.
\]
\[
 \boxed{s=\frac{pr}{2r-p}},\qquad r>\frac p2,\quad s>\frac p2.
\]
$p/2$ 是每一段长度的严格下界，不能由有限焦点弦取到。

\par\medskip
\textbf{\textcolor{CStatement}{与倾斜角的联系}}

若把上方端点叫作 $A$、下方端点叫作 $B$，直线倾斜角为 $0<\theta<\pi$，则
\[
 r=\frac p{1-\cos\theta},\quad
 s=\frac p{1+\cos\theta},\quad
 AB=\frac{2p}{\sin^2\theta}.
\]
交换端点时，两段公式随之交换；倒数和始终不变。

\par\medskip
\textbf{\textcolor{CStatement}{适用范围}}

四种标准开口 $y^2=\pm2px$、$x^2=\pm2py$，若均以正数 $p$ 表示焦准距，长度倒数和都为 $2/p$。
必须是过焦点的有限两点弦；对称轴只有一个交点，不属于此情形。
\end{statementbox}
